\documentclass[11pt]{amsart}

\usepackage[T1]{fontenc}
\usepackage{lmodern}
\usepackage{microtype}
\usepackage{amsmath,amssymb,amsthm}
\usepackage[hidelinks]{hyperref}

\setcounter{MaxMatrixCols}{20}

\hypersetup{
  pdftitle={An Order-14 Counterexample to Conjecture 4 of Pioge et al.}
}

\newtheorem{theorem}{Theorem}
\newtheorem{lemma}[theorem]{Lemma}
\newtheorem{corollary}[theorem]{Corollary}

\newcommand{\HH}{\mathcal H}
\newcommand{\permanent}{\operatorname{per}}
\newcommand{\Real}{\operatorname{Re}}

\title{An Order-14 Counterexample to Conjecture 4 of Pioge et al.}
\subjclass[2020]{15A15, 15A18}
\keywords{permanent, positive semidefinite matrix, Gram matrix, counterexample}
\date{}

\begin{document}

\begin{abstract}
Pioge, Pietrasz, Seron, Novo, and Cerf disproved their Conjecture~4
with a positive semidefinite matrix of order $16$ and reported no
counterexample of order $15$. We give an exact rank-two
Gaussian-integer Gram matrix $A\in\HH_{14}$ such that
\[
  \lambda_{\max}(\Real F_A)>\permanent(A),
  \qquad
  (F_A)_{ij}=a_{ij}\permanent(A(i,j)).
\]
Thus the order of the published construction decreases from $16$ to
$14$; no claim is made for orders at most $13$. The same witness gives
an order-$14$ counterexample to Conjecture~3 of Pioge et al. It is also
diagonally congruent to a rank-two orthogonal projection, yielding
time-delay-induced anomalous bunching with $14$ photons in $14$ modes.
\end{abstract}

\maketitle

\section{The exact counterexample}

Let $\HH_n$ be the cone of $n\times n$ positive semidefinite Hermitian
matrices. For $A=[a_{ij}]\in\HH_n$, let $A(i,j)$ be obtained by
deleting row $i$ and column $j$, and define
\[
  (F_A)_{ij}=a_{ij}\permanent(A(i,j)).
\]
Here $\Real F_A$ denotes the entrywise real part of $F_A$; it is a real
symmetric matrix. Conjecture~4 of Pioge et al.~\cite{PiogeEtAl}
asserts that
\begin{equation}\label{eq:conjecture}
  \lambda_{\max}(\Real F_A)\leq \permanent(A)
  \qquad (A\in\HH_n).
\end{equation}
They disproved \eqref{eq:conjecture} in order $16$ and reported an
unsuccessful numerical search in order $15$. The conjecture is
therefore already known to be false; our contribution is the following
order reduction.

\begin{lemma}\label{lem:rank-two}
Let $X=(x_{rj})\in\mathbb C^{2\times n}$. For $I\subseteq[n]$, let
$X_I$ be the submatrix formed by the columns indexed by $I$, and write
\[
  \prod_{j\in I}(x_{1j}t+x_{2j})
  =\sum_{k=0}^{|I|}q_{I,k}t^k.
\]
If $|I|=|J|=m$, then
\begin{equation}\label{eq:rank-two}
  \permanent(X_I^*X_J)
  =\sum_{k=0}^{m}k!(m-k)!\,\overline{q_{I,k}}q_{J,k}.
\end{equation}
\end{lemma}

\begin{proof}
Expand the two summands in each factor of the permanent. For fixed
$k$ and fixed $k$-subsets $S\subseteq I$ and $T\subseteq J$, exactly
$k!(m-k)!$ bijections from $I$ to $J$ map $S$ onto $T$. Summing first
over such bijections and then over $S$ and $T$ gives
\eqref{eq:rank-two}.
\end{proof}

\begin{theorem}\label{thm:main}
Conjecture~4 of Pioge et al. fails in order $14$, already for a
rank-two Gaussian-integer Gram matrix.
\end{theorem}

\begin{proof}
Let $A=X^*X$, where
$X=X_R+\mathrm iX_I\in\mathbb Z[\mathrm i]^{2\times14}$ and
\begingroup
\scriptsize
\setlength{\arraycolsep}{2.2pt}
\[
X_R=
\begin{pmatrix}
1178&-223&145&278&-388&-78&-9&163&-347&602&-344&529&-168&77\\
-30&-50&68&-14&17&369&866&-686&-467&-88&-133&1199&-794&247
\end{pmatrix},
\]
\[
X_I=
\begin{pmatrix}
317&-59&327&310&426&-109&942&521&-1308&672&-52&117&-497&-368\\
838&744&-92&-307&438&270&-403&-274&-107&907&-728&386&-41&166
\end{pmatrix}.
\]
\endgroup
Then $A\in\HH_{14}$, and
\[
  \det X_{\{1,2\}}=-350880+1045686\,\mathrm i\neq0,
\]
so $\operatorname{rank}A=2$.

Take
\[
  v=(-24,3,24,17,14,-5,13,-16,22,-17,-12,-4,8,-23)^T.
\]
Then $\sum_i v_i=0$ and $v^Tv=3622$. Put $p=\permanent(A)$. Since
no column of $X$ is zero, the polynomial in Lemma~\ref{lem:rank-two}
for $I=[14]$ is nonzero; hence \eqref{eq:rank-two} gives $p>0$.
Applying the lemma to $A$ and to all cofactors $A(i,j)$ gives, in exact
integer arithmetic,
\begin{align*}
\Delta
&:=v^T(\Real F_A)v-p\,v^Tv\\
&=188\,558\,958\,243\,265\,516\,176\,136\,227\,271\,677\,528\,
   258\,899\,031\\
&\qquad{}\cdot10^{45}
 +880\,644\,911\,529\,452\,618\,269\,391\,965\,199\,062\,697\,
  646\,489\,600\\
&>0.
\end{align*}
Therefore
\[
  \lambda_{\max}(\Real F_A)
  \geq \frac{v^T(\Real F_A)v}{v^Tv}
  >p.
\]
For reference, the certified Rayleigh ratio is
\[
  \frac{v^T(\Real F_A)v}{p\,v^Tv}
  =1.001531003891338\ldots.
\]
\end{proof}

\section{Consequences}

\begin{corollary}\label{cor:orders}
Conjecture~4 fails in every order $n\geq14$.
\end{corollary}

\begin{proof}
For $k\geq0$, let $A_k=A\oplus I_k$. Block structure gives
\[
  \permanent(A_k)=p,
  \qquad
  F_{A_k}=F_A\oplus pI_k.
\]
Padding $v$ with $k$ zeros preserves the strict Rayleigh inequality.
\end{proof}

\begin{corollary}\label{cor:conjecture-three}
There is an entrywise positive correlation matrix $B\in\HH_{14}$ such
that
\[
  \permanent(A\circ B)>\permanent(A).
\]
Consequently, Conjecture~3 of Pioge et al.~\cite{PiogeEtAl} also fails
in order $14$.
\end{corollary}

\begin{proof}
Let $\tau=v/\sqrt{3622}$ and, for $\varepsilon>0$, define
\[
  (B_\varepsilon)_{ij}
  =\frac{1+\varepsilon^2\tau_i\tau_j}
  {\sqrt{(1+\varepsilon^2\tau_i^2)
         (1+\varepsilon^2\tau_j^2)}}.
\]
This is the Gram matrix of the unit vectors
\[
  \frac{(1,\varepsilon\tau_i)^T}
       {\sqrt{1+\varepsilon^2\tau_i^2}},
\]
so it is a correlation matrix and is entrywise positive for all
sufficiently small $\varepsilon>0$. Laplace expansion gives every row
and column sum of $F_A$ as $p$. Expanding at $\varepsilon=0$ therefore
yields
\[
  \permanent(A\circ B_\varepsilon)
  =p+\varepsilon^2
  \bigl(\tau^T(\Real F_A)\tau-p\bigr)+O(\varepsilon^4).
\]
The quadratic coefficient is positive by Theorem~\ref{thm:main}.
\end{proof}

\begin{lemma}\label{lem:congruence}
Let $D=\operatorname{diag}(d_1,\ldots,d_n)$ with $d_i>0$, and set
$C=DAD$. Then
\[
  \permanent(C)=\rho\,\permanent(A),
  \qquad
  F_C=\rho F_A,
  \qquad
  \rho=\prod_{i=1}^n d_i^2.
\]
\end{lemma}

\begin{proof}
The first identity follows by extracting $d_i$ from row $i$ and column
$i$. For each $i,j$,
\[
\begin{aligned}
(C)_{ij}\permanent(C(i,j))
&=d_i d_j a_{ij}
  \left(\prod_{r\neq i}d_r\right)
  \left(\prod_{s\neq j}d_s\right)
  \permanent(A(i,j))\\
&=\rho(F_A)_{ij}.
\end{aligned}
\]
\end{proof}

\begin{corollary}\label{cor:projection}
There is a rank-two orthogonal projection $P\in\HH_{14}$ violating
\eqref{eq:conjecture}. It yields time-delay-induced anomalous bunching
with $14$ photons in a $14$-mode interferometer.
\end{corollary}

\begin{proof}
Set
\[
\begin{gathered}
 s=5\,975\,547\,235\,929,
 \quad a=91\,619\,936\,424\,609,\\
 b=170\,572\,929\,732\,992,
 \quad c=119\,336\,587\,939\,769,
\end{gathered}
\]
and
\[
  W=\operatorname{diag}(s,s,s,a,b,s,s,c,s,s,s,s,s,s).
\]
Direct integer multiplication gives
\[
  XWX^*=\eta I_2,
  \qquad
  \eta=144\,241\,568\,073\,276\,041\,742.
\]
Let $D=W^{1/2}$, $Y=\eta^{-1/2}XD$, and
\[
  P=Y^*Y=\eta^{-1}DAD.
\]
Then $YY^*=I_2$, so $P$ is a rank-two orthogonal projection. By
Lemma~\ref{lem:congruence}, $F_P$ and $\permanent(P)$ are obtained from
$F_A$ and $p$ by the same positive scalar factor. Hence the vector
$\tau=v/\sqrt{3622}$ satisfies
\[
  \frac{\tau^T(\Real F_P)\tau}{\permanent(P)}
  =1.001531003891338\ldots>1.
\]

For $d\geq0$, define the Gaussian time-delay Gram matrix
\[
  S_{ij}(d)=\exp\bigl(-d^2(\tau_i-\tau_j)^2\bigr).
\]
Using the row and column sums of $F_P$ and expanding at $d=0$ gives
\begin{align*}
  \frac{\permanent(P\circ S(d))}{\permanent(P)}
  &=1+2\left(
  \frac{\tau^T(\Real F_P)\tau}{\permanent(P)}-1
  \right)d^2+O(d^4)\\
  &=1+0.003062007782676\ldots\,d^2+O(d^4).
\end{align*}
Thus the ratio exceeds $1$ for all sufficiently small $d>0$, which is
temporal anomalous bunching in the criterion of Pioge, Novo, and
Cerf~\cite{PiogeNovoCerf}.

Finally, the two rows of $Y\in\mathbb C^{2\times14}$ are orthonormal
and can be completed to a unitary $U\in\mathbb C^{14\times14}$. For
the first two output modes,
\[
  P_{ij}=\sum_{k=1}^2\overline{U_{ki}}U_{kj}.
\]
Injecting one photon into each input mode therefore realizes the
claimed $14$-photon, $14$-mode configuration. This answers the
fewer-than-$16$-photon question posed in \cite{PiogeNovoCerf}.
\end{proof}

\begin{thebibliography}{9}

\bibitem{PiogeEtAl}
L.~Pioge, K.~K.~Pietrasz, B.~Seron, L.~Novo, and N.~J.~Cerf,
\emph{A logical implication between two conjectures on matrix permanents},
Linear Algebra Appl. \textbf{725} (2025), 309--318.
\href{https://doi.org/10.1016/j.laa.2025.07.011}
{doi:10.1016/j.laa.2025.07.011}.

\bibitem{PiogeNovoCerf}
L.~Pioge, L.~Novo, and N.~J.~Cerf,
\emph{Limits of multimode bunching for boson sampling validation:
anomalous bunching induced by time delays},
arXiv:2601.13792 (2026).
\href{https://arxiv.org/abs/2601.13792}
{arXiv:2601.13792}.

\end{thebibliography}

\end{document}